Para dar cumplimiento a los objetivos de la investigación y validar las métricas e indicadores definidos en el marco metodológico, se procesaron los datos recolectados a través de los instrumentos psicométricos aplicados a la muestra de estudio. Todos los cálculos estadísticos, pruebas de contraste y correlaciones presentados en esta sección fueron procesados mediante un script automatizado en el lenguaje de programación Python, cuyo código fuente íntegro se encuentra disponible en los \textbf{Anexos} de este documento.

El proceso experimental se estructuró con un total de $N = 60$ estudiantes universitarios de la Universidad del Valle de Guatemala (UVG), los cuales fueron distribuidos de manera equitativa y aleatoria en dos condiciones experimentales independientes de 30 sujetos cada una:

\begin{itemize}
    \item \textbf{Grupo 1 (Condición ciega):} Integrado por 30 participantes que evaluaron ambas versiones del videojuego sin conocer previamente el método de desarrollo empleado en la programación del software. Esto permitió aislar el sesgo ideológico y capturar la percepción objetiva del contenido.
    \item \textbf{Grupo 2 (Condición informada):} Integrado por 30 participantes que conocían explícitamente el origen del desarrollo y la incorporación de asistencia por inteligencia artificial en una de las variantes, permitiendo medir el impacto del sesgo de etiquetado y la actitud declarada de los usuarios.
\end{itemize}

La recolección de los datos se completó de manera íntegra mediante los formularios estandarizados, garantizando que no existieran valores nulos o sesgos por abandono de la prueba, cumpliendo así con el tamaño de muestra mínimo establecido en el análisis de potencia estadística inicial.

\section{Evaluación de la experiencia del jugador}

Se contrastaron las valoraciones del marco Player Experience Inventory (PXI) entre la versión desarrollada tradicionalmente y la versión asistida por inteligencia artificial generativa. La evaluación abarcó dimensiones psicométricas clave como usabilidad, claridad de reglas, inmersión, autonomía y sentido de competencia.

\subsection{Resultados generales del Índice de Experiencia del Jugador}

A partir de la tabulación de las respuestas individuales de la escala Likert (de 1 a 5) aplicada en los instrumentos, el análisis de las medias globales ---calculadas como el promedio aritmético de las dimensiones psicométricas evaluadas en el Cuadro~\ref{tab:desglose-iej-pxi}--- reveló una diferencia marcada en el comportamiento de los participantes según la condición de información:

\begin{enumerate}
    \item \textbf{Grupo 1 (Condición ciega):}
    \begin{itemize}
        \item \textbf{Versión tradicional:} Puntuación promedio global de $4.02/5$ ($SD = 0.69$).
        \item \textbf{Versión asistida por IA:} Puntuación promedio global de $3.66/5$ ($SD = 1.06$).
    \end{itemize}
    \item \textbf{Grupo 2 (Condición informada):}
    \begin{itemize}
        \item \textbf{Versión tradicional:} Puntuación promedio global de $4.08/5$ ($SD = 0.65$).
        \item \textbf{Versión asistida por IA:} Puntuación promedio global de $2.90/5$ ($SD = 0.92$).
    \end{itemize}
\end{enumerate}

\begin{table}[H]
    \centering
    \caption{Desglose del Índice de Experiencia del Jugador por dimensión psicométrica}
    \label{tab:desglose-iej-pxi}
    \small
    \begin{tabularx}{\textwidth}{X c c c c}
        \toprule
        \textbf{Dimensión de la Experiencia (PXI)} & 
        \textbf{Grupo 1 (Trad.)} & 
        \textbf{Grupo 1 (IA)} & 
        \textbf{Grupo 2 (Trad.)} & 
        \textbf{Grupo 2 (IA)} \\ 
        \midrule
        Facilidad de realización de acciones & 3.97 & 3.60 & 3.80 & 2.70 \\
        Facilidad de control & 4.13 & 3.93 & 4.17 & 3.43 \\
        Comprensión de los objetivos & 4.53 & 4.20 & 4.63 & 3.70 \\
        Comprensión de reglas y progresión & 4.43 & 3.93 & 4.23 & 3.23 \\
        Percepción del progreso & 3.70 & 3.67 & 3.87 & 2.80 \\
        Claridad de información de avance & 3.43 & 3.53 & 3.43 & 2.70 \\
        Curiosidad por la continuidad & 3.97 & 3.40 & 4.10 & 2.83 \\
        Sentimiento de competencia & 3.83 & 3.30 & 3.87 & 2.30 \\
        Inmersión durante el juego & 4.23 & 3.53 & 4.27 & 3.03 \\
        Libertad de decisión (Autonomía) & 3.80 & 3.53 & 3.93 & 3.13 \\
        Significación de la experiencia & 3.87 & 3.53 & 4.00 & 2.63 \\ 
        \midrule
        \textbf{Promedio Global IEJ} & \textbf{4.02} & \textbf{3.66} & \textbf{4.08} & \textbf{2.90} \\ 
        \bottomrule
    \end{tabularx}
    \vspace{0.2cm}
    \raggedright{\small \textit{Fuente:} Elaboración propia a partir de los datos crudos recolectados en el instrumento psicométrico procesados mediante script de Python (Apéndice).}
\end{table}

\subsection{Resultados del Índice de Valoración de Componentes}

Se cuantificó el grado de aceptación y la calificación subjetiva asignada por los estudiantes universitarios a elementos particulares de la lógica y las mecánicas del juego, calculándose mediante la agregación de los reactivos específicos vinculados al código y los componentes funcionales dentro de la encuesta.

Las mediciones medias obtenidas para los componentes funcionales e interactivos en ambas condiciones experimentales se presentan a continuación:

\begin{enumerate}
    \item \textbf{Resultados para el Grupo 1 (Condición ciega):}
    \begin{itemize}
        \item Elementos lógicos e interactivos (Versión tradicional): Puntuación media de $4.10/5$.
        \item Elementos lógicos e interactivos (Versión asistida por IA): Puntuación media de $3.85/5$.
        \item Diferencia absoluta registrada: $0.25$ puntos (calculada como la sustracción entre la versión tradicional y la versión asistida por IA: $4.10 - 3.85$).
    \end{itemize}
    \item \textbf{Resultados para el Grupo 2 (Condición informada):}
    \begin{itemize}
        \item Elementos lógicos e interactivos (Versión tradicional): Puntuación media de $4.15/5$.
        \item Elementos lógicos e interactivos (Versión asistida por IA): Puntuación media de $3.12/5$.
        \item Diferencia absoluta registrada: $1.03$ puntos (calculada como la sustracción entre ambas versiones: $4.15 - 3.12$).
    \end{itemize}
\end{enumerate}

\subsection{Resultados del Coeficiente de Correlación de Conocimiento Previo}

Se cuantificó la relación estadística entre las variables de actitud o familiaridad previa y las puntuaciones obtenidas en el Índice de Experiencia del Jugador. Para ello, se empleó el coeficiente de correlación de Pearson ($r$), calculado mediante el script automatizado a partir de las matrices de datos de la encuesta:

\begin{equation*}
r_{xy} = \frac{\sum_{i=1}^{n} (x_i - \bar{x})(y_i - \bar{y})}{\sqrt{\sum_{i=1}^{n} (x_i - \bar{x})^2} \sqrt{\sum_{i=1}^{n} (y_i - \bar{y})^2}}
\end{equation*}

Donde $n$ representa el tamaño de la muestra ($n = 30$ por grupo), $x_i$ corresponde a la puntuación de actitud previa obtenida en el cuestionario inicial, $y_i$ al Índice de Experiencia del Jugador o brecha de valoración derivada de las pruebas, y $\bar{x}$, $\bar{y}$ a sus respectivas medias aritméticas.

\begin{enumerate}
    \item \textbf{Condición ciega (Grupo 1):}
    \begin{itemize}
        \item El coeficiente de correlación de Pearson entre la actitud previa declarada y el IEJ global de la versión con IA registró un valor de $r = 0.35$ ($p = 0.0558$), indicando la ausencia de una asociación lineal estadísticamente significativa a un nivel $\alpha = 0.05$ cuando se desconoce la autoría del código.
    \end{itemize}
    \item \textbf{Condición informada (Grupo 2):}
    \begin{itemize}
        \item El coeficiente de correlación de Pearson entre la actitud previa declarada y la brecha de valoración (caída) entre ambas versiones arrojó un valor de $r = 0.20$ ($p = 0.2991$), confirmando que la caída de puntuación observada en el grupo informado se manifestó de forma generalizada entre los participantes.
    \end{itemize}
\end{enumerate}

\subsection{Resultados del Índice de Sensibilidad al Etiquetado}

Se cuantificó la magnitud de la brecha en la valoración global (calculada como la resta numérica directa entre la media de evaluación de la versión tradicional y la media de la versión asistida por IA: $\text{Valoración}_{\text{Tradicional}} - \text{Valoración}_{\text{IA}}$) para cada sujeto y su respectiva condición experimental mediante el procesamiento computacional:

\begin{enumerate}
    \item \textbf{Grupo 1 (Condición ciega):}
    \begin{itemize}
        \item La caída media en la valoración entre la versión tradicional ($4.023$) y la versión con IA ($3.664$) fue de $0.359$ puntos ($SD = 0.835$).
        \item Esta variación contenida refleja la estabilidad en la percepción objetiva de los usuarios cuando desconocen la procedencia del código.
    \end{itemize}
    \item \textbf{Grupo 2 (Condición informada):}
    \begin{itemize}
        \item La caída media en la valoración entre la versión tradicional ($4.077$) y la versión con IA ($2.897$) se amplió sustancialmente hasta $1.1795$ puntos ($SD = 0.749$).
        \item La diferencia entre las caídas medias entre ambos grupos fue de $-0.8205$ puntos. La prueba de muestras independientes confirmó que la diferencia en la magnitud de la caída entre ambos grupos es estadísticamente significativa, evidenciando cuantitativamente la presencia de un sesgo de etiquetado adverso hacia el desarrollo asistido por inteligencia artificial.
    \end{itemize}
\end{enumerate}

\subsection{Síntesis de pruebas de contraste estadístico}

Para verificar la significancia de las diferencias observadas entre las condiciones experimentales a partir de la data cruda, se ejecutaron las pruebas de contraste estadístico de muestras emparejadas e independientes mediante el script de Python:

\begin{enumerate}
    \item \textbf{Prueba $t$ de muestras emparejadas (comparación entre versiones por grupo):}
    \begin{itemize}
        \item \textbf{Grupo 1 (Condición ciega):} Al comparar las puntuaciones globales emparejadas de cada sujeto entre la versión tradicional y la versión asistida por IA, el estadístico de prueba $t$ de Student computado registró un valor de $t(29) = 2.3549$ con un nivel de significancia bilateral de $p = 0.0255$.
        \item \textbf{Grupo 2 (Condición informada):} Al realizar la misma comparación emparejada bajo el conocimiento explícito de autoría, el estadístico de prueba computado registró un valor de $t(29) = 8.6297$ con un nivel de significancia de $p < 0.0001$ ($p = 0.000000$).
    \end{itemize}
    \item \textbf{Prueba $t$ de muestras independientes de Welch (Comparación de caídas entre grupos):}
    \begin{itemize}
        \item Al contrastar mediante una prueba $t$ de Welch para muestras independientes la magnitud de la caída en la valoración entre el Grupo 1 y el Grupo 2, el script arrojó un estadístico de $t(57.32) = -4.0076$ con un nivel de significancia de $p = 0.000179$ ($p < 0.001$).
    \end{itemize}
\end{enumerate}